\chapter{Разработка параллельной OpenMP-версии программного комплекса HeatSim3D}

\section{Анализ информационных зависимостей вычислительного алгоритма схемы FTCS}

\subsection{Независимость вычислений во внутренних узлах на одном шаге}

Расчётный модуль HeatSim3D принимает параметры задачи через аргументы командной строки, выполняет цикл по времени и сохраняет результаты в файлы формата CSV. Каждый шаг по времени состоит из трёх действий. Сначала по формуле~\eqref{eq:ftcs} вычисляются новые значения во всех внутренних узлах, затем на шести гранях восстанавливаются граничные значения, и, наконец, с заданной периодичностью в файлы записываются срезы поля и температура в центре тела.

Первое действие занимает основную часть времени расчёта, так как число внутренних узлов равно $(N_x - 2)(N_y - 2)(N_z - 2)$ и для сетки $150^3$ превышает три миллиона. Как отмечено в первой главе, значение в каждом узле нового слоя зависит только от значений предыдущего слоя. Поэтому между итерациями цикла по внутренним узлам отсутствуют зависимости по данным, и цикл может выполняться параллельно без изменения результата. Между шагами по времени, напротив, существует строгая зависимость, поскольку следующий шаг использует результат предыдущего. Таким образом, распараллеливанию подлежит обработка узлов внутри шага, а сам цикл по времени остаётся последовательным.

\subsection{Организация двух буферов температурного поля в памяти}

Температурное поле хранится в одномерном массиве \texttt{std::vector<double>} длиной $N_x N_y N_z$, а узлу с индексами $(i, j, k)$ соответствует элемент с номером $i N_y N_z + j N_z + k$. При такой нумерации соседние по индексу $k$ узлы расположены в памяти подряд, что позволяет процессору эффективно использовать кэш при последовательном проходе по внутреннему циклу.

Для корректной параллельной работы используются два массива одинакового размера. Первый массив содержит значения на текущем слое и во время шага только читается, а во второй массив записываются новые значения. После завершения шага массивы меняются ролями с помощью функции \texttt{std::swap}, которая обменивает внутренние указатели и не копирует данные. При такой организации каждый элемент нового массива записывается ровно одним потоком, а читаемый массив не изменяется, поэтому гонка данных невозможна. Синхронизация требуется только перед обменом массивов, и её обеспечивает неявный барьер в конце директивы \texttt{parallel for}.

\section{Распараллеливание основных вычислительных циклов расчётного модуля}

\subsection{Слияние пространственных циклов для равномерной загрузки потоков}

В исходной версии программы цикл по внутренним узлам состоял из трёх вложенных циклов по индексам $i$, $j$ и $k$. Если распараллелить только внешний цикл, то число его итераций $N_x - 2$ при небольших сетках сравнимо с числом потоков, и часть потоков может получить на одну итерацию больше других. В версии прошлого семестра для объединения двух внешних циклов использовался параметр \texttt{collapse(2)}. Однако этот параметр появился только в стандарте OpenMP 3.0, а компилятор Microsoft Visual C++ с ключом \texttt{/openmp} поддерживает стандарт 2.0, и программа в этой среде не собиралась бы.

Поэтому в новой версии слияние циклов выполнено вручную. Вводится единый индекс $m$, пробегающий значения от 0 до $(N_x - 2)(N_y - 2) - 1$, а исходные индексы восстанавливаются по формулам
\begin{equation}
    i = 1 + \left\lfloor \frac{m}{N_y - 2} \right\rfloor, \qquad j = 1 + m \bmod (N_y - 2),
\end{equation}
после чего внутри итерации последовательно обрабатывается весь столбец узлов по индексу $k$. Директива \texttt{parallel for} со статическим распределением применяется к циклу по $m$. Уже для сетки $50^3$ число итераций равно 2304, поэтому части, достающиеся потокам, практически одинаковы. Внутренний цикл по $k$ проходит по памяти подряд и хорошо оптимизируется компилятором.

\subsection{Параллельное наложение граничных условий и вычисление средней температуры}

Наложение граничных условий выполняется после каждого шага и состоит из трёх пар циклов, каждая из которых заполняет две противоположные грани куба. Объём этой работы пропорционален $N^2$, а не $N^3$, однако при последовательном выполнении она превращалась бы в постоянную последовательную долю каждого шага. Поэтому каждая пара циклов также снабжена директивой \texttt{parallel for}. Записи в разные элементы массива не пересекаются, и синхронизация внутри этих циклов не нужна.

При каждой записи результатов программа вычисляет среднюю температуру по объёму тела, то есть сумму значений во всех узлах. Для этого цикла использована редукция \texttt{reduction(+:sum)}, при которой каждый поток суммирует свою часть массива, а частичные суммы объединяются по окончании цикла~\cite{chandra2001parallel}. Из-за изменения порядка сложения средняя температура может отличаться от последовательного значения на величину порядка $10^{-13}$~°C. Такое отличие не проявляется в файлах результатов, где значения записываются с четырьмя знаками после запятой.

\subsection{Управление числом потоков и измерение времени расчёта}

Для исследования производительности в расчётный модуль добавлен параметр командной строки \texttt{--threads}, который вызывает функцию \texttt{omp\_set\_num\_threads}. Если параметр не указан, используется число потоков, выбранное средой OpenMP по умолчанию, обычно равное числу ядер процессора. Фактическое число потоков выводится в заголовке протокола расчёта и сохраняется вместе со всеми параметрами задачи в файл \texttt{params.csv}. Это позволяет по файлам результатов однозначно установить, в каком режиме был выполнен расчёт.

Время расчёта измеряется средствами библиотеки \texttt{std::chrono} от начала до конца цикла по времени, поэтому выделение памяти и разбор параметров в него не входят. Измеренное время выводится по окончании расчёта и также записывается в \texttt{params.csv}. При сборке без поддержки OpenMP получается последовательная версия программы, в которой параметр \texttt{--threads} не действует, а в протоколе указывается один поток.

\section{Обеспечение кроссплатформенности программного комплекса для Linux и Windows}

\subsection{Сборка расчётного модуля различными компиляторами с поддержкой OpenMP}

Исходный код расчётного модуля написан на стандартном языке C++17 и содержит единственный платформенно-зависимый фрагмент, связанный с созданием каталога для результатов. В операционной системе Windows для этого используется функция \texttt{\_mkdir}, а в Linux --- функция \texttt{mkdir}, и выбор между ними выполняется директивами условной компиляции. Обращения к библиотеке OpenMP также помещены в блоки условной компиляции, поэтому один и тот же файл собирается как в параллельном, так и в последовательном режиме.

Для Linux расчётный модуль собирается компилятором GCC с ключом \texttt{-fopenmp}. Для Windows исполняемый файл получен кросс-компилятором MinGW-w64 на основе GCC 13 со статической компоновкой, благодаря чему он не требует установки дополнительных библиотек на компьютере пользователя. Сборка в среде Microsoft Visual C++ также возможна после включения поддержки OpenMP в свойствах проекта. Графический интерфейс автоматически выбирает исполняемый файл расчётного модуля в зависимости от операционной системы, а для обеих систем подготовлены сценарии сборки, которые также упаковывают интерфейс в один исполняемый файл.

\subsection{Средства построения графиков и автоматизации вычислительных экспериментов}

Для подготовки результатов к отчёту разработан модуль \texttt{report\_plots.py} на языке Python, использующий библиотеку matplotlib~\cite{bisong2019matplotlib}. Модуль строит графики в едином оформлении, в котором кривые различаются типом линии и маркерами, а шрифт на рисунках совпадает со шрифтом основного текста по начертанию и размеру. Помимо графиков динамики температуры и профилей, модуль вычисляет аналитическое решение для проверки модели и выполняет серию запусков расчётного модуля с разным числом потоков для оценки ускорения.

Все вычислительные эксперименты, результаты которых приведены в третьей главе, выполняются одним сценарием \texttt{experiments.py}. Сценарий запускает расчёты, сохраняет исходные данные замеров в файлы CSV и строит по ним рисунки. Такая организация делает результаты воспроизводимыми, поскольку любой рисунок отчёта можно получить повторно на другом компьютере одной командой. Корректность параллельной версии дополнительно контролируется набором автоматических тестов на основе библиотеки pytest, которые проверяют совпадение результатов при разном числе потоков, работу параметра \texttt{--threads} и соответствие аналитическому решению.
